\documentclass[border=4pt]{standalone}
\usepackage{amsmath,amssymb,mathtools,xcolor,varwidth}
\definecolor{deflectblue}{HTML}{1E6FE0}
\begin{document}
\begin{varwidth}{28em}
\large\bfseries
By \textcolor{deflectblue}{Law I} the body leaving $B$ would run straight on to $c$ (so $Bc$ continues the line $AB$), yet it turns to $C$: some impulse deflected it by $cC$. Now $\triangle SBc=\triangle SAB$ because $Sc$ and $AB$ share the base $SB$ and $Bc$ lies on $AB$; but we are GIVEN $\triangle SBC=\triangle SAB$, hence $\triangle SBC=\triangle SBc$. By \textcolor{deflectblue}{Euclid I.40, equal triangles on the same base SB lie between the same parallels}, so $C$ and $c$ are equidistant from line $SB$ — that is, $cC \parallel BS$. The deflecting impulse (by \textcolor{deflectblue}{Law II}, along $cC$) therefore points exactly along $BS$, straight toward the centre.
\end{varwidth}
\end{document}
